% ===== 2. 関連研究 (Related Work) =====
% 役割:
%   - LLM 応答評価の 3 系統を整理し、本研究の位置付けを明示
%   - 査読で必ず聞かれる「既存手法と何が違うか」への先回り回答
%
% 分量目安: 1 ページ (2 段組換算で約 800 字)
%
% 詳細サーベイ・全比較表は docs/related_metrics_comparison.md を参照
% (本論文の external reference として機能、camera-ready で URL 化)

\section{関連研究}
\label{sec:related_work}

LLM/NLG 応答の自動評価指標は、計算機構の観点から
(a) 表層類似度系、
(b) 埋め込み類似度系、
(c) LLM 駆動評価系
の 3 系統に大別できる。本節では各系統の代表的指標を概観し、本研究が
どの空白を埋めるかを示す。

\subsection{表層・埋め込み類似度指標}
\label{sec:rw:surface_embedding}

BLEU~\cite{papineni2002bleu} および ROUGE~\cite{lin2004rouge} は
n-gram overlap に基づく古典的指標であり、計算コストの低さと再現性から
広く採用されてきたが、構造的整合性と内容的妥当性を単一スコアに混合する
という限界を持つ。BERTScore~\cite{zhang2020bertscore} と
BLEURT~\cite{sellam2020bleurt} は事前学習言語モデルの埋め込み距離または
回帰モデルにより paraphrase 耐性を獲得したが、出力は依然として単一
スカラ (BLEURT) または P/R/F1 の限定的分解 (BERTScore) であり、応答の
どの側面 (構造 / 内容) で失点したかの診断情報を提供しない。

\subsection{LLM 駆動評価}
\label{sec:rw:llm_judge}

近年、評価そのものに LLM を用いる手法が急速に普及している。
G-Eval~\cite{liu2023geval} は GPT-4 に評価基準と Chain-of-Thought
を与え、確率重み付きスコアを生成する。Prometheus~\cite{kim2024prometheus,
kim2024prometheus2} は GPT-4 由来データで open-source LM を fine-tune
し、rubric 条件付きスコアと自然言語フィードバックを返す。
RAGAS~\cite{es2024ragas} は RAG パイプライン特化で 4 軸 (faithfulness,
answer relevancy, context precision/recall) に分解する。MT-Bench を
含む LLM-as-Judge~\cite{zheng2023judging} は GPT-4 を judge として
holistic スコアを生成する。これらは柔軟性が高い反面、(i) 評価軸の
分離が prompt / rubric 設計に完全依存し architectural な保証を持たない、
(ii) judge LLM のバージョン更新で評価が変動する再現性問題、という
共通の限界を抱える。LLM 評価論全体の俯瞰は Chang らのサーベイ
\cite{chang2024survey} に詳しい。

\subsection{本研究の位置付け}
\label{sec:rw:positioning}

上記いずれの系統も、応答の構造完全性 (どの程度、文として壊れていないか)
と命題カバレッジ (どの程度、問いの核心を拾っているか) を**計算機構の
段階で分離して評価する**枠組みを持たない。表層・埋め込み系は両者を単一
スコアに混合し、LLM 駆動系は分離を prompt 設計に委譲するため、軸の
排他性も網羅性も保証されない。本研究は次の 3 点でこの空白を埋める:
(i) 構造完全性 (S 軸) と命題カバレッジ (C 軸) を 2 次元座標 PoR として
architectural に分離する、(ii) 意味損失関数 $L_{\mathrm{sem}}$ を
3 項 (演算子・未知性・因果構造) に分解することで、判定根拠を再利用可能な
診断 primitive として提供する、(iii) judge LLM を計算経路に持たず、
事前学習文埋め込み (SBert) と cascade matcher のみで評価を完結することで
判定の再現性を担保する。
